\begin{theorem}\label{night:P3K.17}
\StatementRecord{night:P3K.17}{false as stated; correction proved}{S6\_06\_BLIND: GPT-6 Astra (OpenAI, ChatGPT Pro mode, effort Pro), recorded re-derivation in the checkpoint dated 18 November 2026}{\cref{night:P3F.22}, \cref{night:P3F.24}, \cref{night:P3F.32}, \cref{night:P3K.7}}
\par\noindent\textbf{Audit correction, 17 November 2026.} The asserted two-fan equality is false. P3K.65 constructs the full negative-to-positive native Cox comparison and gives classes -3k and 0 for the positive r-loop, and -k and 0 for the actual shorter w-loop. At k=1 the original negative fan has classes 25 and 27, respectively. The omitted second outer map has positive-normal determinant winding -2m; the first outer map is the identity. The local charged-cap construction is retained with its full outer marking. The earlier statement and argument below are preserved as the located prior claim.
\medskip

The negative- and positive-diagonal OS49 families have equal
actual return index for every cap order \(m=3k\),
\(k>0\), every signed primitive twist \(n=\pm1\),
and every additional integral delta twist \(d\delta\).
Their index is zero, their oriented return class is zero, and
their actual mapping tori are products as in P3F.24.
The comparison is an oriented bordism carrying the actual
Euler line; it imposes no spin hypothesis on the comparison
bordism.  The endpoint index structures are the prescribed
ones from the original families.
\end{theorem}

\begin{proof}[Preserved earlier argument; scope corrected above]
\begin{enumerate}[label=\arabic*.]
\item \par\medskip\noindent\textbf{The real master space and its actual resolution maps.}\enspace

\item Order the Cox coordinates as \((z_0,z_4,z_1,z_2)\), corresponding
to rays \(u_0,u_4,u_1,u_2\).  Let an auxiliary circle act with
weights \((1,1,-1,-1)\), and let the physical delta circle have
weights \((-1,0,0,+1)\).  These actions commute.  Write
\(a=(z_0,z_4)\), \(c=(z_1,z_2)\), and
\[
       \mu=\|a\|^2-\|c\|^2.
\]
The auxiliary circle is free on \(\mathbb C^4\setminus\{0\}\).
Consequently its quotient is a smooth oriented real seven-manifold.
The residual delta action is semifree.  Indeed every nonzero
two-by-two minor of the matrix with rows
\((1,1,-1,-1)\) and \((-1,0,0,1)\) is \(\pm1\).
On a coordinate support of rank two this proves that the
two-torus has trivial stabilizer.  A rank-one support lies in
one of the two pairs \((z_4,z_1)\), \((z_0,z_2)\),
and the residual circle fixes it pointwise.  Subtracting the
auxiliary lift on these two supports gives normal weights
\((-1,+1)\) in each case.  Thus the circle orbit has the
ordinary rank-three Hopf normal model along its fixed locus,
with no additional finite-quotient strata.
Its positive moment levels are the negative-diagonal small resolution,
and its negative moment levels are the positive-diagonal small
resolution.  Here are the actual maps proving these identifications.
On the positive level use the complex quotient
\[
   \{a\ne0\}\big/\bigl((a,c)\sim(ta,t^{-1}c)\bigr)
       \longrightarrow
   \{([a],ac^t):\operatorname{im}(ac^t)\subset\mathbb Ca\}.
\]
Its target is the total space of \(\mathcal O(-1)^{\oplus2}\)
on \(\mathbb P^1_a\).  Every complex orbit meets a prescribed
positive moment level in exactly one auxiliary-circle orbit:
the required positive real factor solves
\[
 r^2\|a\|^2-r^{-2}\|c\|^2=\tau,
 \qquad
 r^2=\frac{\tau+\sqrt{\tau^2+4\|a\|^2\|c\|^2}}
                 {2\|a\|^2}.
\]
This is smooth for \(a\ne0\) and \(\tau>0\), including
\(c=0\).  The negative-level map is the same formula with
the two pairs interchanged.  Both maps commute with delta.
The positive stability condition excludes \(z_0=z_4=0\),
so its two fan triangles have common edge \(u_1u_2\);
the negative condition has common edge \(u_0u_4\).
Thus these are the specified two fans with their original
coordinate identifications on the common rank-one matrix locus.

\item Restrict to \(|\mu|\le\epsilon\) and a large fixed radius,
and remove a ball of radius \(\rho\) with \(\rho^2<\epsilon\).
The two moment faces are the truncated small resolutions, the
outer face is their common exterior times an interval, and the
new inner face is
\[
             S^7/S^1_{\rm aux}=\mathbb {CP}^3.
\]
This is an actual real bordism relative to the common exterior.
Conjugating the two auxiliary-negative coordinates makes the
auxiliary action scalar-positive; two complex conjugations
preserve the real orientation.  In the resulting homogeneous
coordinates, put
\[
       A=(z_0,\overline z_2),\qquad
       B=(z_4,\overline z_1).
\]
Delta acts by \((A,B)\mapsto(e^{-it}A,B)\).

\par\medskip\noindent\textbf{The circle orbit, its line, and the explicit joint cap.}\enspace

\item The orbit map on this inner face is
\[
 [A,B]\longmapsto
 \left([A],[B],\frac{\|A\|^2}{\|A\|^2+\|B\|^2}\right),
\]
with the irrelevant factor collapsed at either endpoint.
Thus the actual orbit is \(\mathbb {CP}^1*\mathbb {CP}^1\),
identified with \(S^5\subset\mathbb R^3\oplus\mathbb R^3\)
by \((p,q,\phi)\mapsto(\cos\phi\,p,\sin\phi\,q)\).
The two fixed spheres are its two join ends.  On the free orbit
\(S^2_A\times S^2_B\times(0,1)\), the circle line is exactly
\[
           L=\mathcal O(1)_A\boxtimes\mathcal O(-1)_B,
           \qquad x=h_A-h_B.                         \tag{P3K.17.1}
\]
Indeed a nonzero pair determines the relative tensor of its
two tautological lines; delta acts with negative weight on
that tensor.  Taking the associated line with the physical
positive delta character gives the displayed inverse tensor.
This fixes the sign rather than choosing it from the two
normal models independently.

\item The two base charges are therefore \(+1,-1\).  In particular
the line does not extend over either individual conifold cap:
the restriction image of
\(H^2(D^3_A\times S^2_B)\) in the first boundary has no
\(h_A\) term, and the other cap has no \(h_B\) term.
We retain this obstruction and instead join the fixed spheres.

\item Take \(C=S^2_P\times I\), with meridian sphere \(S^2_Q\),
and the line \(\mathcal O(1)_P\boxtimes\mathcal O(1)_Q\)
on \(C\times S^2_Q\).  At its first end use
\[
              (A,B)=(P,\iota Q),
\]
and at its second end use
\[
              (A,B)=(Q,\iota P),                    \tag{P3K.17.2}
\]
where \(\iota\) is the antipodal map of \(\mathbb {CP}^1\).
Both maps pull (P3K.17.1) to \(p+q\).  These are actual
line isomorphisms: the orthogonal-complement map on lines in
\(\mathbb C^2\), followed by its determinant pairing,
identifies \(\iota^*\mathcal O(1)\) complex-linearly with
\(\mathcal O(-1)\); the second end then uses the canonical
interchange of the two tensor factors.  All maps preserve
the physical principal-circle character.

\item The resulting closed five-dimensional orbit correction is the
mapping torus
\[
 Q_{\rm corr}=T_g(S^2\times S^2),\qquad
           g(A,B)=(\iota B,\iota A).
\]
The map \(g\) preserves \(h_A-h_B\) and the product
orientation.  Its line is glued by the specified tensor
isomorphisms, so this statement retains the actual line,
not merely an invariant cohomology class.

\par\medskip\noindent\textbf{The actual parameter gauge on the master space.}\enspace

\item For the closed OS49 parameter \(\lambda=r^m\), with
\(m=3k\) and signed primitive twist \(n=\pm1\), the
extra twist \(d\delta/m\) contributes only a parameter-independent
scalar to the finite-cap line map.  Its source coordinate is unchanged;
the same scalar extends in the Cox lift and across the joint cap.
It therefore introduces no parameter-crossing line class.  The
period-origin calculation in P3F.32 gives the primitive
\(K\)-turn and the angular line winding \(2mn\).
The \(K\)-period on the ordinary Tate coordinates is
\[
      Z\longmapsto a(q)Z,\qquad
      V\longmapsto q^n r^{2mn}b(q)V,
\]
after incorporating the signed traversal (so the bounded units
include the corresponding signed powers), where those units
have parameter-independent logarithms.  Reindexing
the common fan by its integral vertical period removes the
displayed \(q^n\) by the signed integral vertical translation
\(n\).  On the resulting common Cox charts the
bounded map has the explicit lift
\[
 \begin{aligned}
 z_0&\longmapsto\bigl(a(q)b(q)r^{2mn}\bigr)^{-1}z_0,
       &z_1&\longmapsto a(q)z_1,\\
 z_2&\longmapsto b(q)r^{2mn}z_2,
       &z_4&\longmapsto z_4.
 \end{aligned}                                      \tag{P3K.17.3}
\]
The product of its four multipliers is one, so it fixes the
base monomial \(q=z_0z_1z_2z_4\).  Its torus vector is
\(\log a\,(u_1-u_0)+\log(br^{2mn})\,(u_2-u_0)\),
which is exactly the marked \((Z,V)\) action.  Thus its
parameter-dependent factor is one global physical delta
action; it is not two separately chosen gauges at the two
ends of the exceptional geometry.

\item For the exact quotient, label the unit-square flop cells by
\(v\in\mathbb Z^2\).  Their period lattice is
\(\Gamma=\langle(2,1),(-3,-1)\rangle=\mathbb Z^2\),
and \(3(2,1)+2(-3,-1)=(0,1)=K_{\rm trop}\).
It acts freely and transitively on these cells.  Their diagonal
curves are disjoint, since each has the two maximal cones
belonging to its own square.  Choose disjoint equivariant
tubular neighbourhoods and insert one master chart in each.
The quotient of their disjoint union by \(\Gamma\) is
exactly one representative chart with trivial stabilizer.
There is no residual \(K\)-circle on this correction.
The signed parameter return shifts its cell label by
\(n(0,1)=n(3h+2g)\).  Returning to the chosen representative
by that unique period gives (P3K.17.3).  Reversing the
representative convention inverts the one common gauge at
both ends and leaves the comparison unchanged.

\item Fixing \(|r|\), its angular part is
\(e^{2mn i\theta}\) in that physical circle.  The remaining
unit multipliers are independent of \(\theta\); their
chosen logarithms give equivariant isotopies on the Cox
charts.  The positive real rescalings in Section 1 transport
these isotopies to the moment levels and to the inner face.
Explicitly, follow a Cox map by its unique positive auxiliary
rescaling back to the original value of \(\mu\).  Its
implicit derivative is
\(2r^2\|a\|^2+2r^{-2}\|c\|^2>0\) away from the
removed origin, so this is smooth through the real bordism.
The bounded units can be changed to their values at \(q=0\)
using \(a(tq),b(tq)\), \(0\le t\le1\), with an
exterior cutoff.  The maps stay invertible and commute with
both circles; they are independent of the parameter angle.
This change introduces no angular character on the inner face.
These maps need not preserve a fixed Euclidean radius.
Use transported truncations: under the moment-preserving
equivariant isotopy \(G_t\), take the inner and outer
faces to be \(G_t\) of the original faces.  They stay
embedded and disjoint, while the moment faces stay at
\(\mu=\pm\epsilon\).  Glue these transported charts
and collars by their exact maps.  For the remaining
constant-unit clutch \(G_0\), its chosen logarithms give
\(G_0^t\), and
\((x,t)\mapsto(G_0^t x,e^{2\pi it})\) trivializes
the mapping torus.  The angular delta factor preserves
the round face and is trivial on its circle orbit.
Consequently the two maps (P3K.17.2) have the same angular
gauge, which extends by that constant circle multiplication
over all of \(C\times S^2_Q\).  No difference of two
angular gauges is introduced on its ends.  The closed
correction over the parameter circle therefore carries
a line pulled back from \(Q_{\rm corr}\), after this
explicit gauge identification.  Its first Pontryagin class
is pulled back from the same five-dimensional factor.
The potential mixed line class can also be checked directly.
Cover the correction by its free-orbit cylinder and joint-cap
cylinder; their overlap has two components.  The actual
transition functions there are
\(\varphi_1(\theta)=g(\theta)\varphi_{1,0}\) and
\(\varphi_2(\theta)=g(\theta)\varphi_{2,0}\), with
the same scalar from (P3K.17.3).  Change the cap frame by
\(g(\theta)^{-1}\).  Both angular factors disappear
simultaneously, leaving parameter-independent transitions.
This explicitly excludes a line class on
\(S^1_{\rm joint}\times S^1_{\rm parameter}\).
Both six-dimensional characteristic numbers vanish:
\[
  \langle x^3,[Q_{\rm corr}\times S^1]\rangle=0,
  \qquad
  \langle p_1x,[Q_{\rm corr}\times S^1]\rangle=0.       \tag{P3K.17.4}
\]

\par\medskip\noindent\textbf{The additional normal-sphere boundary in the surgery bordism.}\enspace

\item The fixed locus in the master bordism must also be retained.
There are exactly two relevant auxiliary-quotient fixed charts:
\((z_4,z_1)\) and \((z_0,z_2)\).  Each has auxiliary
weights \((1,-1)\), and its quotient is \(\mathbb R^3\)
by the explicit invariant map
\[
      (z_+,z_-)\longmapsto
       \bigl(2\Re(z_+z_-),2\Im(z_+z_-),
                         |z_+|^2-|z_-|^2\bigr).
\]
The inverse orbit is determined by
\(|z_\pm|^2=(\sqrt{\mu^2+4|z_+z_-|^2}\pm\mu)/2\).
The two moment faces in each chart are discs.  Removing the
inner master-space ball removes an interior three-ball in
each of these fixed charts.  Away from the flop curve the
fixed sphere minus its two endpoint discs is unchanged.
Gluing this annulus times the moment interval proves that
the fixed three-dimensional cobordism is exactly
\[
                 W_F=(S^2\times I)\setminus
                             (\mathring D^3\sqcup\mathring D^3).
\]
The Cox reindexing preserves the two charts; it does not
permute or identify their boundary spheres.

\item The oriented normal three-plane bundle is trivial over this
three-dimensional cobordism, with its boundary product
structures retained.  On the two Cox fixed charts its
positive-weight complex normal bundle is respectively
\(L^{\oplus2}\) and \((L^{-1})^{\oplus2}\), so the
associated three-plane bundle has \(w_2=0\).  The other
two boundary spheres have the original \(SU(2)\) normal
model.  Restriction of \(H^2(W_F;\mathbb Z/2)\) to
the four boundary spheres is injective, so \(w_2\) vanishes
globally.  An oriented three-plane bundle on a three-dimensional
complex with this class zero is trivial: the obstruction on
two-cells is \(\pi_1(SO(3))=\mathbb Z/2\), and that
on three-cells is \(\pi_2(SO(3))=0\).  Boundary
trivializations differ by maps \(S^2\to SO(3)\), all
null-homotopic.  Collar homotopies therefore match the chosen
product trivializations at every boundary.  The actual circle
line restricted to the constant normal section has boundary
values \(0,0,+h_A,-h_B\).

\item Cap the two original-family sphere boundaries by \(D^3\),
as in their actual conifold construction, and join the two
inner-face spheres by the charged cap \(C\) above.  This
closes \(W_F\) to \(S^2\times S^1\).  Crossing with
the parameter circle closes the fixed four-dimensional base
to \(\bar X\).  The described local trivialization makes
it the product with that parameter circle, retaining the
possible parameter-independent clutch along the first
\(S^1\).  In particular it is a mapping torus over the
first circle of \(S^2\), crossed with the parameter circle.
All gluings preserve the sphere orientation.  Its signature
is zero, and its tangent first Pontryagin number is zero.
For the exact orientation comparison, capping the first two
boundaries leaves \(S^2\times I\).  Its line has degrees
\(+1\) and \(-1\) in the two complex boundary markings.
Its product transport consequently has degree \(-1\)
between those markings.  The joint cap has exactly that
antipodal base change at its second end.  The composite has
degree \(+1\), giving the oriented sphere mapping torus
above.  On the meridian sphere the first end map is also
antipodal, so both meridian orientations give the same
physical Euler value one.

\item Apply the explicit collar construction of P3K.7 to the
complement of this fixed locus in the master-space orbit
family.  At the two original-family boundaries it gives
their actual conifold mapping tori.  At the inner face it
gives \(Q_{\rm corr}\times S^1\).  It has one further
boundary, the normal sphere bundle \(B=S(\nu_{\bar X})\).
The meridian-one line is glued by (P3K.17.2) and (P3K.17.3)
through every collar, so it extends across this oriented
bordism.  This accounts for the boundary which would be
missed by capping the two inner fixed spheres separately.

\item The actual parameter-dependent Cox action is delta, which
acts trivially on the circle orbit and its normal rank-three
bundle.  Thus on the closed fixed base both its normal data
and the line data are pulled back from the three-dimensional
\(S^2\)-mapping-torus factor.  With the notation of P3K.7,
\(e_{\mathrm{vert}}=2x-\pi^*b\), it follows that
\(b^2=0\).  The sphere-bundle calculation, valid as an
oriented characteristic calculation, gives
\[
       \frac{4\langle x^3,[B]\rangle
                  -\langle p_1(TB)x,[B]\rangle}{24}
       =\frac{3\langle b^2,[\bar X]\rangle
                    -\langle p_1(T\bar X),[\bar X]\rangle}{24}
       =0.                                             \tag{P3K.17.5}
\]
No value of \(p_1(\nu_{\bar X})\) is assumed here.

\item For a closed oriented six-manifold with a complex line,
\(4x^3-p_1x\) is an oriented characteristic number and
vanishes on oriented boundaries carrying that line.  Hence
the two original-family return values are equal after
(P3K.17.4) and (P3K.17.5).  The actual endpoint manifolds
already have their prescribed spin structures, so their
integer \(A\) is the same characteristic expression.
The signs of the two additional boundary components do not
affect their zero contributions; the two moment faces have
opposite boundary orientations.  Together with P3F.32 this identifies the negative-diagonal
return index with zero for all the stated orders and twists.
The exact Hopf transfer P3F.22 gives its oriented zero class,
and P3F.24 gives the product mapping torus.

\end{enumerate}
\end{proof}